% ============================================================
%  ELECTRÓNICA — conexiones del ESP32 (esquemático, paso 8)
%  Sin librería calc: solo coordenadas explícitas.
% ============================================================
\begin{tikzpicture}[font=\scriptsize, line cap=round, line join=round]
\colorlet{bus}{gray!65}

% ================= CELDA DE MEDICIÓN (izquierda) =================
\draw[dashed, gray!70] (0.3,3.3) rectangle (3.9,7.2);
\node[font=\tiny, text=gray!60] at (2.1,7.05) {celda de medición};
\draw[fill=yellow!30] (0.5,5.9) rectangle (3.7,6.7);
\node at (2.1,6.3) {MH-410D (CO$_2$)};
\draw[fill=yellow!30] (0.5,4.9) rectangle (3.7,5.7);
\node at (2.1,5.3) {MH-441D (CH$_4$)};
\draw[fill=yellow!30] (0.5,3.9) rectangle (3.7,4.7);
\node at (2.1,4.3) {LuminOx (O$_2$)};

% ================= SENSORES AMBIENTALES (arriba) =================
\draw[fill=green!15] (3.4,7.6) rectangle (6.8,8.6);
\node at (5.1,8.1) {2$\times$ SHT45 --- T y HR};
\draw[fill=purple!15] (7.1,7.6) rectangle (10.7,8.6);
\node at (8.9,8.1) {3$\times$ DS18B20 --- T cama};
\draw[fill=brown!15] (11.0,7.6) rectangle (14.2,8.6);
\node at (12.6,8.1) {2$\times$ humedad sustrato};

% ================= ESP32 (centro) =================
\draw[line width=1.2pt, fill=blue!12] (7.0,4.0) rectangle (10.0,6.0);
\node[font=\scriptsize\bfseries] at (8.5,5.35) {ESP32};
\node[font=\tiny, align=center] at (8.5,4.6) {lee sensores, comanda relés,\\registra y envía por USB};

% ================= RELÉS (derecha) =================
\draw[line width=1.1pt, fill=red!8] (11.6,3.6) rectangle (15.4,6.4);
\node[font=\scriptsize\bfseries] at (13.5,6.12) {módulo de 5 relés};
\node[font=\tiny, anchor=north west, align=left] at (11.75,5.85)
  {GPIO32 $\rightarrow$ válvula 6 (aireación)\\
   GPIO33 $\rightarrow$ válvula 9 (3 vías)\\
   GPIO27 $\rightarrow$ válvula 11 (cierre)\\
   GPIO14 $\rightarrow$ PTC (calefacción)\\
   GPIO13 $\rightarrow$ nebulizador};

% ================= PESAJE (abajo) =================
\draw[fill=orange!15] (0.4,1.6) rectangle (4.1,2.4);
\node[align=center] at (2.25,2.0) {4$\times$ celdas de carga (esquinas)\\[-1pt]
  {\tiny puente Wheatstone}};
\draw[fill=orange!30] (4.6,1.6) rectangle (7.4,2.4);
\node[align=center] at (6.0,2.0) {HX711 --- ADC 24 bits};

% ================= ALIMENTACIÓN =================
\draw[fill=gray!10] (0.1,0.05) rectangle (4.7,1.0);
\node[align=center] at (2.4,0.53) {5 V/5 A lógica $\cdot$ 12 V/10 A potencia\\[-1pt]
  {\tiny (PTC y válvulas) --- tierra común}};

% ================= RASPBERRY PI =================
\draw[fill=gray!12] (11.4,1.4) rectangle (15.6,2.2);
\node at (13.5,1.8) {PC del laboratorio --- dashboard};

% ================= CABLES =================
% gas: celda -> ESP32 (3 UART)
\draw[bus, line width=1pt] (3.7,6.3) -- (7.0,5.7);
\node[font=\tiny, anchor=south] at (4.4,6.38) {UART1 $\cdot$ 16/17};
\draw[bus, line width=1pt] (3.7,5.3) -- (7.0,5.3);
\node[font=\tiny, anchor=south] at (4.4,5.34) {UART2 $\cdot$ 25/26};
\draw[bus, line width=1pt] (3.7,4.3) -- (7.0,4.55);
\node[font=\tiny, anchor=north] at (4.4,4.22) {UART0 $\cdot$ 1/3};

% ambientales: arriba -> ESP32
\draw[bus, line width=1pt] (5.1,7.6) -- (5.1,6.9) -- (7.6,6.9) -- (7.6,6.0);
\node[font=\tiny, anchor=south] at (6.35,6.94) {I$^2$C $\cdot$ 21/22};
\draw[bus, line width=1pt] (8.9,7.6) -- (8.9,6.0);
\node[font=\tiny, anchor=east] at (8.78,7.35) {1-Wire $\cdot$ GPIO4};
\draw[bus, line width=1pt] (12.6,7.6) -- (12.6,6.9) -- (9.6,6.9) -- (9.6,6.0);
\node[font=\tiny, anchor=south] at (11.1,6.94) {ADC $\cdot$ 34/35};

% relés: ESP32 -> módulo
\draw[bus, line width=1pt] (10.0,5.0) -- (11.6,5.0);
\node[font=\tiny, anchor=south] at (10.8,5.04) {5$\times$ digital};

% pesaje: celdas -> HX711 -> ESP32
\draw[bus, line width=1pt, -{Stealth[length=2mm]}] (4.1,2.0) -- (4.6,2.0);
\draw[bus, line width=1pt] (6.0,2.4) -- (6.0,3.2) -- (7.8,3.2) -- (7.8,4.0);
\node[font=\tiny, anchor=south] at (6.9,3.24) {GPIO18/19};

% RPi: ESP32 -> RPi (USB)
\draw[bus, line width=1pt] (9.2,4.0) -- (9.2,2.9) -- (13.5,2.9) -- (13.5,2.2);
\node[font=\tiny, anchor=south] at (11.35,2.94) {WiFi / USB-serial};

\end{tikzpicture}
